\documentclass[11pt]{article}
\usepackage{../../../tex/lecturestyle}

\title{Bostan--Mori's Algorithm (Advanced)}
\author{}
\date{}

\begin{document}
\thispagestyle{fancy}

\lecturetitleband{Polynomials \& Formal Power Series}{Bostan--Mori's Algorithm (Advanced)}{8}

\section{Overview}
The companion lecture on the $K$-th term of a linearly recurrent sequence
showed how Bostan--Mori's algorithm finds a single coefficient
$[x^K]\,P(x)/Q(x)$ in $O(M(d)\log K)$. This lecture extends that
machinery two ways. First, a variant --- the \emph{MSB-first} Bostan--Mori
algorithm --- extracts a whole \emph{interval} of coefficients
$[x^{[L,R]}]\,P(x)/Q(x)$ in about the same time, and this one primitive
turns out to solve several other problems too: shifting a linearly
recurrent sequence's generating function, and computing
$x^K \bmod \Gamma(x)$ (the object at the heart of Fiduccia's algorithm) by
an entirely different route. Second, when fast (FFT-based) polynomial
multiplication is available, keeping intermediate polynomials in their
DFT representation across Bostan--Mori's iterations --- rather than
converting back and forth every round --- cuts the constant factor from
$2M(d)\log K$ to about $\tfrac23 M(d)\log K$.

Assumes the companion lecture on the $K$-th term algorithm in full, and
(for the second half) the definition of the discrete Fourier transform
and its use for fast polynomial multiplication.

\section{Setup and Notation}
As before, $A$ is linearly recurrent with characteristic data
$\Gamma(x)=x^d-c_1x^{d-1}-\cdots-c_d$ and
$Q(x)=1-c_1x-\cdots-c_dx^d$, and $A(x)=P(x)/Q(x)$ for some $P$ of degree
$<d$. Ring operations are counted as the unit of work, and $M(d)$ is the
cost of multiplying two degree-$d$ polynomials.

\begin{definitionbox}[(Interval coefficient extraction)]
For a formal power series (or Laurent series) $f(x)=\sum_i f_i x^i$
--- allowing $i<0$, with $f_i:=0$ there by convention --- and integers
$L\le R$, define
\[
  [x^{[L,R]}]f(x) \;:=\; \sum_{i=L}^{R} f_i\, x^{\,i-L},
\]
an ordinary polynomial of degree $\le R-L$, \emph{reindexed to start at
$x^0$}. This reindexing matters: once extracted, the result is treated as
a fresh polynomial with its own $0$-based coefficients, not as a
``window'' that remembers the original degree $L$ it came from.
\end{definitionbox}

\section{Interval Coefficient Extraction}

\subsection{Reducing to a pure inverse}
\begin{definitionbox}
Given $P$ of degree $\le d-1$, $Q$ of degree $\le d$ with $Q(0)=1$, and
$0 \le L \le R$, compute $[x^{[L,R]}]\,P(x)/Q(x)$.
\end{definitionbox}

Since $P$ has degree $<d$, $[x^{[L,R]}](P(x)/Q(x))$ only ever draws on
$1/Q(x)$'s coefficients from degree $L-d+1$ up to $R$ --- lower-degree
terms of $1/Q(x)$ cannot reach degree $L$ once multiplied by $P$'s
degree-$\ge 0$ terms, and terms above $R$ obviously don't contribute.
So it suffices to solve the pure-inverse version:

\begin{definitionbox}
Given $Q$ of degree $\le d$ with $Q(0)=1$, and integers $L \le R$ with
$R \ge 0$ ($L$ possibly negative), compute $[x^{[L,R]}]\,1/Q(x)$.
\end{definitionbox}

\subsection{The MSB-first recursion}
The derivation reuses the even/odd-part identity from the ordinary
Bostan--Mori algorithm: $Q(x)Q(-x) = V(x^2)$ for some $V$ of degree
$\le d$. Then
\[
  \frac{1}{Q(x)} \;=\; \frac{Q(-x)}{Q(x)Q(-x)} \;=\; Q(-x)\cdot\frac{1}{V(x^2)}.
\]
Let $S(x) := [x^{[L-d,\,R]}]\,\dfrac{1}{V(x^2)}$ (an ordinary,
reindexed polynomial, per the definition above). Since $Q(-x)$ has degree
$\le d$, one can check directly from the convolution formula for
$1/Q(x)=Q(-x)\cdot(1/V(x^2))$ that
\[
  [x^{[L,R]}]\,\frac{1}{Q(x)} \;=\; [x^{[d,\,R-L+d]}]\bigl(Q(-x)\cdot S(x)\bigr),
\]
where the right-hand extraction reads $Q(-x)\cdot S(x)$ as an ordinary
product of two already-reindexed polynomials (literal $0$-based
coefficient indices, not a further degree-shifted read --- this is the
step where an implementation most easily goes wrong).

It remains to find $S(x)$, an interval of $1/V(x^2)$ --- a series with
only even-degree terms. Since $[x^{2i}](1/V(x^2))=[x^i](1/V(x))$, setting
\[
  L' = \left\lceil \tfrac{L-d}{2} \right\rceil, \qquad
  R' = \left\lfloor \tfrac{R}{2} \right\rfloor
\]
(chosen so that $L-d\le 2i\le R \iff L'\le i\le R'$) and
$W(x):=[x^{[L',R']}]\,1/V(x)$ recovers $S$ termwise: writing
$j_0 = 2L'-L+d$, $S$'s only nonzero entries are
$S_{j_0+2t} = W_t$ for $t=0,\ldots,|W|-1$. So finding $S$ reduces to
finding $W$ --- the \emph{same kind of problem}, on $V$ instead of $Q$,
with interval length roughly halved.

\begin{theorembox}
Interval extraction $[x^{[L,R]}]\,1/Q(x)$ reduces, in $O(M(d))$ work per
level, to the same problem with $R$ replaced by $\lfloor R/2\rfloor$.
Iterating $O(\log R)$ times bottoms out at $R=0$, where
$[x^{[L,0]}]\,1/Q(x)$ is immediate: since $Q(0)=1$, only degree $0$ can be
nonzero (every degree in $[L,-1]$ is negative, hence $0$ by convention),
so the answer is the single monomial $x^{-L}$ (a $1$ at output position
$-L$). Overall this costs $O(M(d)\log R + M(R-L+d))$, or with more care
about multiplying polynomials of very different degrees,
$O\!\left(M(d)\log R + \tfrac{R-L+1}{d}M(d)\right)$.
\end{theorembox}

\begin{notebox}
The multiplication $Q(-x)\cdot S(x)$ in each level only needs $S$'s
\emph{middle} $R-L+1$ coefficients of the full product (indices $d$
through $R-L+d$) --- a \emph{middle product}, which can be computed in
the same asymptotic cost as computing the two factors' degrees would
suggest, via the transpose principle (see below), rather than computing
the whole product and discarding the rest.
\end{notebox}

% FIGURE: bm_msb
\begin{figure}[H]
\centering
\begin{tikzpicture}[x=1cm,y=1cm]
\node[lnode,minimum width=3.4cm,minimum height=7mm] (n0) at (0,-0.00) {$Q_0$,\; $[L,R]=[5,13]$};
\node[lnode,minimum width=3.4cm,minimum height=7mm] (n1) at (0,-1.45) {$Q_1$,\; $[L,R]=[2,6]$};
\node[lnode,minimum width=3.4cm,minimum height=7mm] (n2) at (0,-2.90) {$Q_2$,\; $[L,R]=[0,3]$};
\node[lnode,minimum width=3.4cm,minimum height=7mm] (n3) at (0,-4.35) {$Q_3$,\; $[L,R]=[-1,1]$};
\node[lnode,minimum width=3.4cm,minimum height=7mm,lhi2] (n4) at (0,-5.80) {$Q_4$,\; $[L,R]=[-1,0]$\\[-1pt]\scriptsize base case: $R=0$, answer $x^{1}$};
\draw[larrow] ([xshift=-1.1cm]n0.south) -- ([xshift=-1.1cm]n1.north);
\draw[larrow,lectureaccent] ([xshift=1.1cm]n1.north) -- ([xshift=1.1cm]n0.south);
\draw[larrow] ([xshift=-1.1cm]n1.south) -- ([xshift=-1.1cm]n2.north);
\draw[larrow,lectureaccent] ([xshift=1.1cm]n2.north) -- ([xshift=1.1cm]n1.south);
\draw[larrow] ([xshift=-1.1cm]n2.south) -- ([xshift=-1.1cm]n3.north);
\draw[larrow,lectureaccent] ([xshift=1.1cm]n3.north) -- ([xshift=1.1cm]n2.south);
\draw[larrow] ([xshift=-1.1cm]n3.south) -- ([xshift=-1.1cm]n4.north);
\draw[larrow,lectureaccent] ([xshift=1.1cm]n4.north) -- ([xshift=1.1cm]n3.south);
\node[llabel,anchor=east,align=right] at (-2.0,-2.9) {descend: $Q_{t+1}(x^2)=Q_t(x)Q_t(-x)$,\\ $L'=\lceil (L-d)/2\rceil,\ R'=\lfloor R/2\rfloor$};
\node[llabel,anchor=west,align=left,text=lectureaccent] at (2.0,-2.9) {unwind: interleave $W$ into $S$,\\ then take the middle product\\ of $Q_t(-x)$ and $S$};
\end{tikzpicture}
\caption{The MSB-first recursion for $[x^{[L,R]}]\,1/Q(x)$ with $d=2$ and $[L,R]=[5,13]$. Each level squares away the odd part of the denominator and roughly halves the interval ($L'=\lceil (L-d)/2\rceil$, $R'=\lfloor R/2\rfloor$) until $R=0$; the answer is then rebuilt on the way back up, one middle product per level.}\label{fig:bm_msb}
\end{figure}

\subsection{This is the transpose of the ordinary algorithm}
\begin{notebox}
Write the target answer as $[x^{[L,R]}]\,1/Q(x) = \sum_i a_i x^i$
(0-indexed, $i=0,\ldots,R-L$), and for a test polynomial $B(x)=\sum_i
b_i x^i$ of the same degree, note that
\[
  [x^R]\,\frac{B(x)}{Q(x)} \;=\; \sum_{i=0}^{R-L} a_i\, b_{R-L-i}.
\]
So ``output $[x^{[L,R]}]\,1/Q(x)$'' and ``given $B(x)$, output
$[x^R]\,B(x)/Q(x)$'' are transpose problems (up to reversing coefficient
order), and by the transpose principle they can be solved with the same
asymptotic arithmetic cost --- and an algorithm for either can be
mechanically converted into one for the other. The second problem is
exactly ordinary (LSB-first) Bostan--Mori. This is also the origin of the
names: writing $1/Q(x) = \prod_{k\ge 0} Q_k(x^{2^k})$ (the infinite
product the doubling identity generates), ordinary Bostan--Mori processes
this product from $k=0$ upward (LSB-first), while the algorithm above
processes it from the top down (MSB-first).
\end{notebox}

\subsection{A byproduct: fast formal power series inverse}
Taking $[L,R]=[0,d]$ recovers the ordinary power series inverse
$1/Q(x) \bmod x^{d+1}$ --- usually computed via Newton's method in
$O(M(d))$. A version of the algorithm above, adapted so that both $L<0$
intervals reduce to $[0,R]$ intervals and $Q$'s coefficients above degree
$R$ are simply dropped, halves the working size every level and achieves
the same $O(M(d))$ bound by this route instead.

\section{Two Applications}

\subsection{Shifting a linearly recurrent sequence}
\begin{definitionbox}
Given $P(x)$, $Q(x)$, and $K\ge 0$, find a rational generating function
for $(A_K, A_{K+1}, A_{K+2}, \ldots)$.
\end{definitionbox}

The shifted sequence $B_i := A_{K+i}$ satisfies the \emph{same}
recurrence as $A$ (it's still generated by unrolling the same $c_1,
\ldots, c_d$), so $B(x) = S(x)/Q(x)$ for some $S$ of degree $<d$. To find
$S$: compute $[x^{[0,d-1]}]B(x) = [x^{[K,K+d-1]}]\,P(x)/Q(x)$ (an
interval extraction of the \emph{original} $A(x)=P(x)/Q(x)$), multiply by
$Q(x)$, and keep the terms of degree $<d$ --- this recovers exactly $S$,
since $S(x) = B(x)Q(x) \bmod x^d$ by definition of $B=S/Q$. Total cost:
$O(M(d)\log K)$.

\subsection{\texorpdfstring{$x^K \bmod \Gamma(x)$}{x\^{}K mod Gamma(x)} via Bostan--Mori}
\begin{definitionbox}
Given $K \ge 0$, compute $x^K \bmod \Gamma(x)$.
\end{definitionbox}

Fiduccia's algorithm already solves this by binary exponentiation modulo
$\Gamma(x)$. Here is a second route, through interval extraction. Let
$U(x) := [x^{[K-d+1,K]}]\,1/Q(x)$, and let $R(x) = \sum_i r_i x^i$ be the
target $x^K \bmod \Gamma(x)$, with $\widetilde{R}(x)$ its coefficients
reversed. Both $U$ and $\widetilde R$ turn out to relate to $A_K$: the
companion lecture's Fiduccia theorem gives $A_K = [x^{d-1}](A(x)\,
\widetilde R(x))$, while $A_K = [x^K]P(x)/Q(x) = [x^{d-1}](P(x)\,U(x))$
follows directly from $A_K=[x^K]P(x)\cdot(1/Q(x))$ and $P$'s degree bound.
Since $P(x)=A(x)Q(x)$, substituting gives
$[x^{d-1}](A(x)\widetilde R(x)) = [x^{d-1}](A(x)Q(x)U(x))$ for
\emph{every} choice of initial values $A_0,\ldots,A_{d-1}$ (as $\widetilde
R, Q, U$ don't depend on them) --- which forces the polynomial identity
\[
  \widetilde R(x) \;=\; Q(x)\,U(x) \bmod x^d .
\]
So: compute $U$ by interval extraction, multiply by $Q$, truncate to
degree $<d$, and reverse the coefficients. Cost: $O(M(d)\log K)$, the
same as Fiduccia, and (as with the $K$-th-term case) known to be smaller
by a constant factor in practice.

\section{Constant-Factor Speedup via DFT Representations}
The remaining idea applies to the \emph{ordinary} (single-coefficient)
Bostan--Mori algorithm when fast polynomial multiplication is available.
Each iteration computes $U(x)=P(x)Q(-x)$ and $V(x)=Q(x)Q(-x)$, naively
two full multiplications, i.e.\ about $2M(d)$ work via $3$ length-$2n$
DFT/IDFT pairs each (a product of two degree-$d$ polynomials by FFT costs
roughly $3$ transforms of size $2n\ge 2d{+}1$: two forward, one
inverse) --- so $6$ transforms per iteration if done naively twice over.
Keeping everything in DFT form throughout avoids re-transforming what's
already been transformed:

\begin{itemize}[leftmargin=1.5em]
  \item $Q(-x)$'s length-$2n$ DFT is just $Q(x)$'s length-$2n$ DFT with
        indices shifted by $n$ (since evaluating at $-\zeta^i=\zeta^{i+n}$
        for a primitive $2n$-th root $\zeta$) --- free, no transform
        needed.
  \item Pointwise multiplication in DFT space gives $F_{2n}[U]$ and
        $F_{2n}[V]$ directly from $F_{2n}[P]$, $F_{2n}[Q]$, and the
        shifted $F_{2n}[Q]$ --- no transform needed.
  \item Extracting the even (or odd) part of a polynomial from its
        length-$2n$ DFT, as a length-$n$ DFT of that part, is an $O(n)$
        butterfly-style combination of paired entries --- no full
        transform needed.
  \item Only going back \emph{up} in size --- from a length-$n$ DFT to a
        length-$2n$ DFT of the same polynomial, needed to prepare next
        iteration's inputs --- genuinely costs transforms: one length-$n$
        IDFT (recover the polynomial) and one length-$n$ DFT evaluated at
        a rotated set of points (get the odd-indexed entries), for each
        of $P'$ and $Q'$.
\end{itemize}

Tallying only the genuinely necessary transforms --- $4$ of length $n$
per iteration, versus $6$ of length $2n$ for the naive two-multiplication
approach --- brings each iteration's cost down to about
$\tfrac23 M(d)$, and the whole algorithm to about
$\tfrac23 M(d)\log K$.

\begin{notebox}
This is a specific, well-known instance of a broader and very useful
pattern: when a computation is a pipeline of polynomial operations linked
by multiplications, look for opportunities to keep intermediate results
in DFT form and only transform back when a non-multiplicative operation
(here, taking an even/odd part) genuinely requires it. The same idea
extends, with more bookkeeping, to the interval-extraction algorithm of
earlier sections.
\end{notebox}

\section{Reference Implementations}
Both examples below implement interval coefficient extraction
(cross-checked against brute-force power series inversion, including
negative $L$), the shift application, and $x^K \bmod \Gamma(x)$ via
Bostan--Mori (cross-checked against Fiduccia's binary exponentiation).
The DFT-representation speedup of Section 5 is a constant-factor,
implementation-heavy optimization built on top of an actual FFT/NTT
engine; it is described above with its complexity analysis but not
implemented here, consistent with how the source material itself treats
it as a refinement layered on the core algorithm.

\codenote
\subsection{Rust}
\corelisting{Rust}{example.rs}

\subsection{Python}
\corelisting{Python}{example.py}

\section{Worked Example}
\begin{examplebox}
For $Q(x)=1-x-x^2$ (the Fibonacci recurrence) and any $P$, both
$x^K \bmod \Gamma(x)$ computed via interval extraction and via Fiduccia's
binary exponentiation agree exactly across hundreds of random $(Q,K)$
pairs with $K$ up to $2000$ in the reference implementations, as does the
shift application: computing $S(x)$ so that $S(x)/Q(x)$ generates
$(A_K,A_{K+1},\ldots)$ and unrolling it directly reproduces the same
values as unrolling the original sequence from index $K$ onward.
\end{examplebox}

\section{Further Reading}
\begin{itemize}[leftmargin=1.5em]
  \item AtCoder Algorithm Lectures index (source of this curriculum):
        \url{https://info.atcoder.jp/entry/algorithm_lectures/index}
  \item A. Bostan, R. Mori, ``A Simple and Fast Algorithm for Computing
        the N-th Term of a Linearly Recurrent Sequence'', SOSA 2021,
        pp.\ 118--132.
  \item A. Sch\"onhage, ``Variations on computing reciprocals of power
        series'', Information Processing Letters 74(1--2), 2000,
        pp.\ 41--46.
\end{itemize}

\end{document}
